%% kapa: exam=1 round=32 year=2021 period=2 form=A subject=accounting subject_ko=회계학 seq=18 number=58
%% source: Q-Net 감정평가사 기출문제 (cst00309) artlSeq=5208397 file=extracted/1008/5208397/A형/2021년도 제32회 감정평가사 1차 2교시 A형.pdf
%% conversion: kordoc 4.24.51 -> tools/gen1.py (auto); review=proofread-sonnet; flags=none
\begin{question}{58}
(주)감평이 20×1년 말 재무상태표에 계상하여야 할 충당부채는? (단, 아래에서 제시된 금액은 모두 신뢰성 있게 측정되었다.)

{\small\centering\begin{tabular}{|p{0.44\linewidth}|p{0.44\linewidth}|}\hline
\multicolumn{1}{|c|}{사건} & \multicolumn{1}{c|}{비고} \\ \hline
20×1년 9월 25일에 구조조정 계획이 수립되었으며 예상비용은 ￦300,000으로 추정된다. & 20×1년 말까지는 구조조정계획의 이행에 착수하지 않았다. \\ \hline
20×1년 말 현재 소송이 제기되어 있으며, 동 소송에서 패소 시 배상하여야 할 손해배상금액은 ￦200,000으로 추정된다. & (주)감평의 자문 법무법인에 의하면 손해발생 가능성은 높지 않다. \\ \hline
미래의 예상 영업손실이 ￦450,000으로 추정된다. & \\ \hline
회사가 사용 중인 공장 구축물 철거 시, 구축물이 정착되어 있던 토지는 원상복구의무가 있다. 원상복구원가는 ￦200,000으로 추정되며 그 현재가치는 ￦120,000이다. & \\ \hline
판매한 제품에서 제조상 결함이 발견되어 보증비용 ￦350,000이 예상되며, 그 지출가능성이 높다. 동 보증은 확신유형 보증에 해당한다. & 예상비용을 보험사에 청구하여 50\% 만큼 변제받기로 하였다. \\ \hline
\end{tabular}\par}
\vspace{4pt}
\begin{choices}
  \item ￦295,000
  \item ￦470,000
  \item ￦550,000
  \item ￦670,000
  \item ￦920,000
\end{choices}
\end{question}
\answer{2}
